%%% SECTION 9: A CANCER PATIENT'S JOURNEY THROUGH PHYSICAL AI %%%

A Cancer Patient's Journey, Physical AI \cite{patient-journey-paper} provides
definitive evidence that Physical AI can manage a complete oncology trial
patient journey autonomously across ten stages, with documented safety
outcomes, regulatory compliance, and cost savings. This single-patient
simulation documents a Stage IIIB non-small cell lung cancer (NSCLC) patient
over 1,120 days, demonstrating end-to-end Physical AI capability from
prescreening through closeout. Together with the 24-hour multi-patient
simulation \cite{site-docs}, this constitutes pivotal evidence that
state-of-the-art AI has the ability to understand, work with, and scale
Physical AI applications beyond human and prior technology capabilities.

\subsection{Ten-Stage Pipeline Architecture}
\label{subsec:journey-pipeline}

The patient journey is organized into ten sequential stages, each documented
with regulatory mapping, clinical activities, Physical AI system involvement,
and outcomes. The stages span the complete trial lifecycle, demonstrating that
Physical AI is not limited to isolated robotic procedures but can orchestrate
the entire patient experience.

\begin{table}[H]
\centering
\caption{Ten-Stage Physical AI Patient Journey Pipeline}
\label{tab:journey-stages}
\small
\begin{tabularx}{\textwidth}{c l X}
\toprule
\textbf{Stage} & \textbf{Name} & \textbf{Physical AI Activities} \\
\midrule
1 & Prescreening & AI-assisted eligibility assessment, record review \\
2 & Screening & Physical AI diagnostic imaging, biomarker analysis \\
3 & Informed Consent & AI-enhanced consent with robot demonstrations \\
4 & Enrollment & Automated registration, baseline data collection \\
5 & Treatment Planning & Digital twin modeling, treatment optimization \\
6 & Surgical Intervention & Robot-assisted tumor resection \\
7 & Post-Surgical Care & Robotic monitoring, rehabilitation support \\
8 & Follow-up Treatment & AI-guided therapy, companion robot support \\
9 & Long-term Monitoring & Automated surveillance, imaging follow-up \\
10 & Closeout & Automated reporting, regulatory documentation \\
\bottomrule
\end{tabularx}
\end{table}

The ten-stage pipeline demonstrates end-to-end capability that differentiates
Physical AI oncology trials from prior robotic surgery studies focusing only on
surgical outcomes. Each stage is documented with three-perspective ASCII
diagrams showing the regulatory, clinical, and patient views simultaneously.

\subsection{Pre-Trial Stages: Prescreening Through Enrollment}
\label{subsec:journey-pretrial}

Stages 1 through 4 cover the pre-trial process from initial patient
identification through formal enrollment. Stage 1 (Prescreening) uses AI
algorithms to review patient medical records against trial eligibility
criteria, identifying candidates with Stage IIIB NSCLC who meet inclusion
criteria and have no disqualifying conditions. Physical AI systems assist by
automating record retrieval, cross-referencing imaging databases, and flagging
potential candidates for physician review.

Stage 2 (Screening) deploys Physical AI imaging systems for diagnostic
confirmation, including robotic CT-guided biopsy for tissue sampling and
AI-assisted pathology analysis for tumor characterization. Stage 3 (Informed
Consent) implements the Physical AI-specific consent requirements from the
adapted 21 CFR Part 50 \cite{cfr50-adapt}, including robot demonstrations,
AI transparency explanations, and documentation of patient rights under
AB 2847 from the site documentation package \cite{site-docs}.

Stage 4 (Enrollment) uses automated systems for patient registration, baseline
data collection, and treatment arm assignment. The regulatory framework mapping
shows how each pre-trial stage complies with the adapted standards: the
Adaption: ICH Harmonised Guideline \cite{ich-adapt} governs investigator
conduct at each stage, the adapted 21 CFR Part 50 governs informed consent and
human subjects protection, and the adapted 21 CFR Part 312 \cite{cfr312-adapt}
governs IND compliance and Phase 0 validation completion.

\subsection{Treatment Stages: Surgical and Therapeutic Interventions}
\label{subsec:journey-treatment}

Stages 5 through 8 cover the active treatment period. Stage 5 (Treatment
Planning) uses digital twin technology to create a patient-specific
computational model of the tumor and surrounding anatomy, enabling simulation
of multiple treatment approaches before selecting the optimal strategy. AI
algorithms analyze the digital twin against historical outcomes data to predict
treatment response.

Stage 6 (Surgical Intervention) documents the robot-assisted tumor resection
performed by a surgical robot system. The procedure is monitored through the
Physical AI safety matrix requirements, with continuous force, position, and
vital sign monitoring. Surgical outcomes include procedure duration, blood loss,
margin clearance, and immediate post-operative status. The adverse event
profile documents any robot-related events during the procedure.

Stage 7 (Post-Surgical Care) deploys rehabilitation exoskeletons for mobility
support, companion robots for emotional support, and monitoring cobots for
continuous vital sign surveillance. Stage 8 (Follow-up Treatment) covers
adjuvant therapy with AI-guided dosing optimization and Physical AI-assisted
treatment delivery.

USL scoring \cite{usl-standard} is applied during treatment stages to verify
that each robot involved in patient care meets minimum readiness thresholds.
Surgical robots must achieve Advanced band (USL 7.0 or higher) for Dimension D
before performing procedures, while support robots (cobots, companions) may
operate at Intermediate band (USL 5.0 or higher).

\subsection{Post-Treatment Stages: Closeout and Analysis}
\label{subsec:journey-closeout}

Stages 9 and 10 cover the post-treatment period through trial closeout. Stage 9
(Long-term Monitoring) uses automated Physical AI surveillance systems for
scheduled imaging, blood work monitoring, and symptom assessment over the
1,120-day follow-up period. AI algorithms analyze trends in monitoring data to
detect potential recurrence or late treatment effects. Stage 10 (Closeout)
automates the generation of regulatory documentation including case report
forms, safety narratives, and protocol compliance reports, ensuring that all
Physical AI-generated data is properly archived per the retention requirements
established in the adapted standards.

\subsection{Regulatory Compliance Coverage}
\label{subsec:journey-compliance}

The patient journey demonstrates compliance with all three adapted regulatory
standards at every stage. The regulatory framework mapping documents which
specific provisions of each adapted standard apply at each stage, creating a
traceability matrix from clinical activities to regulatory requirements.

\begin{table}[H]
\centering
\caption{Regulatory Standard Coverage Across Journey Stages}
\label{tab:journey-compliance}
\small
\begin{tabularx}{\textwidth}{c X X X}
\toprule
\textbf{Stage} & \textbf{ICH E6(R3)} & \textbf{21 CFR Part 50} & \textbf{21 CFR Part 312} \\
\midrule
1-2 & Investigator qualifications & Subject identification & IND eligibility \\
3-4 & Informed consent procedures & Consent elements & Enrollment procedures \\
5-6 & Quality management & Safety matrix & Phase reporting \\
7-8 & Safety reporting & Ongoing consent & Adverse event reporting \\
9-10 & Record keeping & Data protection & Annual/final reports \\
\bottomrule
\end{tabularx}
\end{table}

\subsection{Physical AI Ecosystem and Data Flow}
\label{subsec:journey-ecosystem}

Data flows between Physical AI systems, clinical staff, regulatory systems, and
patients throughout the ten-stage journey. The Physical AI ecosystem comprises
multiple robot types operating in coordinated workflows, with data exchange
managed through MCP servers \cite{national-mcp-paper} using the Model Context
Protocol \cite{mcp-protocol}. Each robot generates data streams that are
captured, validated, and routed to appropriate destinations: clinical databases
for care coordination, regulatory systems for compliance reporting, and
research databases for outcome analysis.

The data flow architecture ensures that no single robot operates in isolation.
Surgical robots share procedure data with monitoring cobots, which in turn
update patient records accessed by companion robots for post-procedure support.
This integrated data flow supports the continuous monitoring requirements of
the adapted 21 CFR Part 50 \cite{cfr50-adapt} and the data governance
provisions of the adapted ICH E6(R3) \cite{ich-adapt}.

\subsection{FDA Cost-Savings Analysis}
\label{subsec:journey-cost}

The patient journey includes a detailed financial analysis documenting
per-patient cost reductions at each stage. Physical AI reduces costs through
automated screening (Stage 1-2), streamlined consent processes (Stage 3),
efficient enrollment (Stage 4), AI-optimized treatment planning (Stage 5),
precision surgical intervention reducing complications and hospital stay
(Stage 6), automated post-surgical monitoring reducing nursing workload
(Stage 7), AI-guided follow-up reducing unnecessary clinic visits (Stage 8-9),
and automated closeout documentation reducing administrative burden (Stage 10).

Tufts CSDD data \cite{tufts-delay} establishes that each day of delay in drug
development has significant economic consequences. The timeline compression
demonstrated in the patient journey, from faster screening and enrollment
through more efficient treatment delivery and automated documentation,
translates directly to accelerated drug development timelines. NCI data on
trial costs \cite{nci-trials-paying} provides context for understanding the
economic burden that Physical AI can reduce.

\subsection{Trial Timeline Comparison}
\label{subsec:journey-timeline}

The Physical AI trial timeline introduces the Phase 0 simulation validation
stage but achieves net time savings through compression at every subsequent
phase. Enrollment speed increases through automated screening and AI-assisted
eligibility verification. Treatment throughput increases through 24-hour
robotic operations and reduced per-procedure times. Data processing accelerates
through automated documentation and real-time regulatory reporting through MCP
servers. The net effect is a significantly shorter overall trial timeline
despite the addition of the Phase 0 requirement.

\subsection{Discussion of Patient Journey Evidence}
\label{subsec:journey-discussion}

The patient journey evidence is significant for several reasons. First, it
demonstrates autonomous workflow capability across all ten stages, establishing
that Physical AI is not limited to individual procedures but can manage
complete trial operations. Second, it validates the regulatory adaptations by
showing that a real patient scenario can be mapped to the adapted standards at
every stage. Third, it provides the financial data that underpins the economic
analysis in Section~\ref{sec:financial}. Fourth, it establishes safety
performance benchmarks against which future Physical AI trials can be measured.

Both the single-patient journey and the 24-hour multi-patient simulation are
pivotal completed examples demonstrating state-of-the-art AI ability to
understand, proficiently work with, and scale Physical AI applications. The
evidence is comprehensive: safety data, regulatory compliance, cost analysis,
and patient outcomes all support the transition to Physical AI oncology trials.

\subsection{Limitations and Future Directions}
\label{subsec:journey-limitations}

The single-patient simulation has acknowledged limitations: it documents a
simulated rather than real patient, covers only one cancer type (Stage IIIB
NSCLC), and represents outcomes from AI-generated scenarios rather than
clinical observations. These limitations are partially addressed by the
24-hour multi-patient simulation covering 15 cancer types and 168 patients.
Full validation requires real-world clinical trials, which this National
Platform is designed to enable through its comprehensive regulatory, standards,
infrastructure, and documentation framework.
